\hypertarget{r15.-h-rogers_j-from-the-rank-one-mean-value-theorem-l1-the-siegel-rogers-formula-for-delta-independent-j-frames}{%
\subsection*{\texorpdfstring{R15. (H-Rogers)\(_j\) from the rank-one mean value theorem (L1): the Siegel-Rogers formula for \(\Delta\)-independent \(j\)-frames}{R15. (H-Rogers)\_j from the rank-one mean value theorem (L1): the Siegel-Rogers formula for \textbackslash Delta-independent j-frames}}\label{r15.-h-rogers_j-from-the-rank-one-mean-value-theorem-l1-the-siegel-rogers-formula-for-delta-independent-j-frames}}

\begin{quote}
\textbf{STATUS.} Theorem R below (= (H-Rogers)\(_j\) of R14 for \(1\le j<k\), constant exactly 1, all Borel \(F\ge0\), identity in \([0,\infty]\)) is \textbf{PROVED {[}P{]}} from: (L1) in ranks \(k,k-1,\dots,k-j+1\) (all \(\ge2\)) for a \emph{coherent} family of groups (axioms (G1)-(G2), verified for the concrete families in 0.3), the hypothesis \textbf{(h1\(^\ell\)): every nonzero LEFT ideal of \(\mathcal O\) is principal}, and standard Haar-measure facts (cited {[}L{]}, from memory, classical). Written in one pass under a 22 minute limit, NOT refereed. Two bookkeeping caveats are isolated and treated in §4: the left/right form of (h1) (Lemma 4.1a: for a unimodular \(\Lambda_W\), e.g. Leech, (h1) \(\iff\) (h1\(^\ell\))), and the exact definition of the norm-one group in {[}G{]} when a real place exists (Lemma 4.2a; irrelevant for \(\mathbb K=\mathbb H\)). No numerical check was made (none is natural here). No transitivity of \(SL_k(\mathcal O)\) on primitive vectors, no zeta function and no volume computation is used: the constant is fixed orbit by orbit by (L1) in rank \(k\) itself.
\end{quote}

Tags: {[}P{]} proved here; {[}L{]} cited (source stated); {[}GAP{]} not proved. {[}G{]} = Gargava arXiv:2107.04844 (NOT opened in this session; its Theorem 16 is quoted from \texttt{../R6\_quaternionic/proof\_revised.md} Thm 4.1/Rem. 4.1a and \texttt{../R14\_leech\_tensor/novelty\_check.md} item 4d). {[}R14{]} = \texttt{../R14\_leech\_tensor/proof.md}.

\hypertarget{setting-sides-groups}{%
\subsubsection*{0. Setting, sides, groups}\label{setting-sides-groups}}

\textbf{0.1 Data.} \(\Delta\) a finite-dimensional division algebra over \(\mathbb Q\), \(\Delta_{\mathbb R}=\Delta\otimes\mathbb R\cong\mathbb K_1\times\dots\times\mathbb K_s\), \(\mathbb K_v\in\{\mathbb R,\mathbb C,\mathbb H\}\) (as in {[}R14{]}), \(\mathcal O\subset\Delta\) an order. Row vectors: \(\Delta_{\mathbb R}^{1\times m}\); scalars act on the LEFT, matrices \(\gamma\in GL_m(\Delta_{\mathbb R})=\prod_vGL_m(\mathbb K_v)\) on the RIGHT (\(y\mapsto y\gamma\)); the two actions commute. \(\mathcal O^{1\times m}\) is a free left \(\mathcal O\)-module. Lebesgue measure on \(\Delta_{\mathbb R}^{a\times b}\) is always normalised by \(\mathrm{covol}(\mathcal O^{a\times b})=1\); these normalisations are compatible with products (\(\mathcal O^{a\times b}=\prod\mathcal O\)). \(\mathrm{rk}=\mathrm{rk}_\Delta\) is the rank of a matrix over the division ring \(\Delta\) (left row rank = right column rank {[}L, classical{]}).

\textbf{(h1\(^\ell\))} every nonzero left ideal of \(\mathcal O\) is principal (\(=\mathcal Oc\)). (See 4.1 for the relation with the right-ideal form (h1) of {[}R14{]}.)

\textbf{0.2 Lemma (norm) {[}P{]}.} For \(a\in GL_m(\mathbb K)\) put \(\|a\|:=|\det_{\mathbb R}(x\mapsto xa\ \text{on}\ \mathbb K^{1\times m})|\). Then \(\|a\|=|\det_{\mathbb R}(x\mapsto ax\ \text{on}\ \mathbb K^{m\times1})|\). For \(\gamma=(\gamma_v)\in GL_m(\Delta_{\mathbb R})\) put \(\|\gamma\|=\prod_v\|\gamma_v\|\); then right multiplication by \(\gamma\) on \(\Delta_{\mathbb R}^{a\times m}\) and left multiplication on \(\Delta_{\mathbb R}^{m\times b}\) multiply Lebesgue measure by \(\|\gamma\|^a\), \(\|\gamma\|^b\). If \(\delta\in GL_m(\mathcal O)\) then \(\|\delta\|=1\).
\emph{Proof.} Both expressions are continuous homomorphisms \(GL_m(\mathbb K)\to\mathbb R_{>0}\), trivial on elementary (unipotent) matrices and equal to \(\prod|c_i|^{\dim\mathbb K}\) on \(\mathrm{diag}(c_1,\dots,c_m)\); \(GL_m\) of a division ring is generated by elementary and diagonal matrices (Gauss elimination). \(\delta\in GL_m(\mathcal O)\) maps the lattice \(\mathcal O^{1\times m}\) onto itself. \(\square\)

\textbf{0.3 Coherent families of groups.} A \emph{coherent family} is a sequence of closed subgroups \(\mathsf G_m\subseteq GL_m(\Delta_{\mathbb R})\), \(m\ge1\), with

\begin{itemize}
\tightlist
\item
  \textbf{(G1)} \(\mathsf G_m\) is normal in \(GL_m(\Delta_{\mathbb R})\) and \(\|\gamma\|=1\) for \(\gamma\in\mathsf G_m\);
\item
  \textbf{(G2)} for \(m\ge2\): \(\{\gamma\in\mathsf G_m:e_1\gamma=e_1\}=\mathsf H^{(m)}:=\Big\{u(n,h):=\begin{pmatrix}1&0\\ n&h\end{pmatrix}:\ n\in\Delta_{\mathbb R}^{(m-1)\times1},\ h\in\mathsf G_{m-1}\Big\}\).
\end{itemize}

\(\Gamma_m:=\mathsf G_m\cap GL_m(\mathcal O)\) (discrete). Examples {[}P{]}: if \(\chi_m:GL_m(\Delta_{\mathbb R})\to T\) are continuous homomorphisms to abelian groups with \(\chi_m(u(n,h))=\chi_{m-1}(h)\) and \(\ker\chi_m\subseteq\{\|\cdot\|=1\}\), then \(\mathsf G_m=\ker\chi_m\) is coherent (\(e_1\gamma=e_1\) iff \(\gamma=u(n,h)\), \(h\in GL_{m-1}\)). This covers

\begin{enumerate}
\def\labelenumi{(\Alph{enumi})}
\tightlist
\item
  \(\mathsf G^1_m=\prod_vSL_m(\mathbb K_v)\) (reduced norm one at every place; \(\chi=(\mathrm{nrd}_v)_v\), multiplicative on block-triangular matrices);
\item
  \(\mathsf G^{+}_m=\{\gamma:\det_{\mathbb R}(R_\gamma|\Delta_{\mathbb R}^{1\times m})=1\}\) and \(\mathsf G^{\pm}_m=\{\|\gamma\|=1\}\) (the real determinant of \(x=(x_1|x'')\mapsto(x_1+x''n\,|\,x''h)\) is \(\det R_h\)).
  If no \(\mathbb K_v\) is \(\mathbb R\) (Hurwitz, icosians, Leech: \(\mathbb K=\mathbb H\)), \(GL_m(\mathbb K_v)\) is connected, real determinants are positive, and \(\mathsf G^+_m=\mathsf G^\pm_m=\{N=1\}\) for every usual definition of the real norm \(N\) (a positive power of \(\|\cdot\|\)): this is the group \(\mathsf G_k\) of {[}R14{]} and (as recorded in {[}R6, Rem. 4.1a(i){]}) of {[}G{]}. Family (A) is in general a smaller group (for \(s\ge2\) or \(\mathbb K=\mathbb C\)). \textbf{The theorem below holds for every coherent family for which (L1) holds}; which one is meant in {[}G{]} could not be re-read here (see 4.2).
\end{enumerate}

\textbf{0.4 The random lattice; sides (point (a)).} \(\mathcal O^{1\times k}\gamma\) depends only on the coset \(\Gamma_k\gamma\), so the space of lattices is \(X_k:=\Gamma_k\backslash\mathsf G_k\) (the "\(\mathsf G_k/\Gamma_k\)" of {[}R14{]} must be read this way). Transposition is an anti-isomorphism \(M_k(\Delta)\to M_k(\Delta^{op})\), \((y\gamma)^T=\gamma^T\cdot_{op}y^T\); it maps \(\Gamma_k\backslash\mathsf G_k\) to \(\mathsf G_k^{op}/\Gamma_k^{op}\), rows \(\mathcal O^{1\times k}\gamma\) to columns \(\gamma^T(\mathcal O^{op})^k\), and \(R_\gamma\) on rows is the same real-linear map as \(L_{\gamma^T}\) on columns. So the family \(\{\mathcal O^{1\times k}\gamma\}\) of {[}R14{]} IS Gargava\textquotesingle s family \(\{g\mathcal O'^k\}\) for \((\Delta^{op},\mathcal O^{op})\) {[}P, given that both use the same norm condition{]}.

\textbf{(L1)\(_m\) {[}L; {[}G, Thm 16 + Prop. 14{]} for \((\Delta^{op},\mathcal O^{op})\), read from the source in rounds R6/R14, not in this session; proof not re-verified{]}.} \(m\ge2\). \(\Gamma_m\) is a lattice in \(\mathsf G_m\); let \(\mu_m\) be the invariant probability measure on \(X_m\). For \(f\in C_c(\Delta_{\mathbb R}^{1\times m})\): \(\int_{X_m}\sum_{0\ne y\in\mathcal O^{1\times m}}f(y\gamma)\,d\mu_m=\int f(x)\,dx\).

\textbf{0.5 Lemma (Borel extension) {[}P{]}.} (L1)\(_m\) implies the same identity in \([0,\infty]\) for every Borel \(f:\Delta_{\mathbb R}^{1\times m}\to[0,\infty]\).
\emph{Proof.} \(\lambda(A):=\int_{X_m}\#\{y\ne0:y\gamma\in A\}d\mu_m\) is a Borel measure (monotone convergence; the integrand is a countable sum of Borel functions), with \(\int f\,d\lambda=\int f\,dx\) for \(f\in C_c\); so \(\lambda\) is Radon and equals Lebesgue measure (Riesz). Then \(\int_{X_m}\sum f(y\gamma)=\int f\,d\lambda\) for Borel \(f\ge0\) by monotone convergence. \(\square\)

\hypertarget{statement}{%
\subsubsection*{1. Statement}\label{statement}}

\textbf{Theorem R {[}P, from (L1)\(_k\),...,(L1)\(_{k-j+1}\) and (h1\(^\ell\)){]}.} Let \((\mathsf G_m)\) be a coherent family, \(1\le j<k\), and assume (h1\(^\ell\)) and (L1)\(_m\) for \(m=k,k-1,\dots,k-j+1\). Then for every Borel \(F:\Delta_{\mathbb R}^{j\times k}\to[0,\infty]\)
\[\int_{\Gamma_k\backslash\mathsf G_k}\ \sum_{Y\in\mathcal O^{j\times k},\ \mathrm{rk}_\Delta Y=j}F(Y\gamma)\ d\mu_k(\gamma)=\int_{\Delta_{\mathbb R}^{j\times k}}F(y)\,dy\qquad(\text{in }[0,\infty]).\tag{R$_{j,k}$}\]
In particular (H-Rogers)\(_j\) of {[}R14{]} holds for \(2\le j<k\) (there \(k>r\ge j\), so all ranks used are \(\ge k-r+1\ge2\)).

Induction on \(j\). (R\(_{1,k}\)) is (L1)\(_k\) + Lemma 0.5. Induction step: for \(j\ge2\) (so \(k\ge3\)), (R\(_{j-1,k-1}\)) + (L1)\(_k\) + {[}\(\Gamma_{k-1}\) is a lattice in \(\mathsf G_{k-1}\){]} \(\Rightarrow\) (R\(_{j,k}\)); and (R\(_{j-1,k-1}\)) needs (L1) in ranks \(k-1,\dots,k-j+1\). Not circular: (L1)\(_k\) is used as a black box for the first row only.

\hypertarget{lemmas}{%
\subsubsection*{2. Lemmas}\label{lemmas}}

\textbf{Lemma 2.1 (first row, uses (h1\(^\ell\))) {[}P{]}.} Every \(0\ne y\in\mathcal O^{1\times k}\) is \(y=a\,e_1\delta\) with \(a\in\mathcal O\setminus\{0\}\), \(\delta\in GL_k(\mathcal O)\). ("\(y'=e_1\delta\) primitive": \(\Delta y'\cap\mathcal O^{1\times k}=\mathcal Oy'\).)
\emph{Proof.} \(I:=\{c\in\Delta:cy\in\mathcal O^{1\times k}\}\supseteq\mathcal O\) is a left \(\mathcal O\)-submodule of \(\Delta\), finitely generated (\(Iy\) is a subgroup of \(\mathcal O^{1\times k}\)), so \(tI\) is a left ideal for some integer \(t>0\), \(tI=\mathcal Oc_0\), and \(M:=\Delta y\cap\mathcal O^{1\times k}=Iy=\mathcal Oy'\) with \(y'=t^{-1}c_0y\); \(y=ay'\), \(a\in\mathcal O\). \(Q:=\mathcal O^{1\times k}/M\) is a finitely generated torsion-free left \(\mathcal O\)-module (\(M\) is saturated), hence free by the left-handed version of {[}R14, Lemma 1.1{]} (same proof with sides exchanged, using (h1\(^\ell\))); so \(\mathcal O^{1\times k}=\mathcal Oy'\oplus N\) with \(N\) free, of rank \(k-1\) (tensor with \(\mathbb Q\)). Let \(\delta\in M_k(\mathcal O)\) have rows \(y'\) and a basis of \(N\). \(x\mapsto x\delta\) is a left-\(\mathcal O\)-linear bijection of \(\mathcal O^{1\times k}\); its inverse is right multiplication by a matrix in \(M_k(\mathcal O)\). So \(\delta\in GL_k(\mathcal O)\), \(e_1\delta=y'\). \(\square\)

\textbf{Lemma 2.2 (the stabiliser) {[}P{]}.} Let \(k\ge3\), \(\mathsf H=\mathsf H^{(k)}\) as in (G2), and assume \(\Gamma_{k-1}\) is a lattice in \(\mathsf G_{k-1}\) (hence \(\mathsf G_{k-1}\) is unimodular {[}L, classical{]}); \(dh\) := Haar measure of \(\mathsf G_{k-1}\) with \(\mathrm{vol}(\Gamma_{k-1}\backslash\mathsf G_{k-1})=1\), \(dn\) := normalised Lebesgue measure on \(V=\Delta_{\mathbb R}^{(k-1)\times1}\). Then: (i) \(u(n,h)u(n',h')=u(n+hn',hh')\), i.e. \(\mathsf H\cong\mathsf G_{k-1}\ltimes V\); (ii) \(d\eta:=dn\,dh\) is a left and right Haar measure on \(\mathsf H\) (\(\mathsf H\) unimodular); (iii) \(\Gamma_{\mathsf H}:=\mathsf H\cap\Gamma_k=\{u(n,h):n\in\mathcal O^{(k-1)\times1},h\in\Gamma_{k-1}\}\); (iv) if \(\mathcal F'\) is a Borel fundamental domain for \(\Gamma_{k-1}\backslash\mathsf G_{k-1}\) and \(\mathcal P\) a fundamental parallelotope of \(\mathcal O^{(k-1)\times1}\) in \(V\), then \(\{u(n,h):h\in\mathcal F',n\in\mathcal P\}\) is a fundamental domain for \(\Gamma_{\mathsf H}\backslash\mathsf H\); so \(\eta(\Gamma_{\mathsf H}\backslash\mathsf H)=1\).
\emph{Proof.} (i) matrix product. (ii) Left translation by \(u(n_0,h_0)\) is \((n,h)\mapsto(n_0+h_0n,h_0h)\); \(n\mapsto h_0n\) preserves \(dn\) by (G1) and Lemma 0.2. Right translation is \((n,h)\mapsto(n+hn_0,hh_0)\): a translation in \(n\) for fixed \(h\), and \(dh\) is right invariant. (iii) \(u(n,h)^{-1}=u(-h^{-1}n,h^{-1})\); so \(u(n,h)\in GL_k(\mathcal O)\) iff \(n\) integral and \(h\in GL_{k-1}(\mathcal O)\). (iv) \(u(m,\sigma)u(n,h)=u(m+\sigma n,\sigma h)\): there is a unique \(\sigma\in\Gamma_{k-1}\) with \(\sigma h\in\mathcal F'\) and then a unique \(m\in\mathcal O^{(k-1)\times1}\) with \(m+\sigma n\in\mathcal P\). \(\square\)

\textbf{Lemma 2.3 (torus lemma) {[}P{]}.} Let \(A\in\mathcal O^{p\times q}\) with \(\mathrm{rk}_\Delta A=p\), and \(\psi:\Delta_{\mathbb R}^{p\times1}\to[0,\infty]\) Borel, \(\mathcal P\) a fundamental parallelotope of \(\mathcal O^{q\times1}\). Then \(\int_{\mathcal P}\sum_{b\in\mathcal O^{p\times1}}\psi(b+An)\,dn=\int_{\Delta_{\mathbb R}^{p\times1}}\psi(t)\,dt\).
\emph{Proof.} \(\tilde\psi(t):=\sum_b\psi(b+t)\) is a Borel function on \(\mathbb T'=\Delta_{\mathbb R}^{p}/\mathcal O^{p}\) with \(\int_{\mathbb T'}\tilde\psi=\int\psi\) (Tonelli). \(\alpha:\mathbb T=\Delta_{\mathbb R}^{q}/\mathcal O^{q}\to\mathbb T'\), \(n\mapsto An\), is a well defined continuous homomorphism (\(A\mathcal O^q\subseteq\mathcal O^p\)), and it is onto: column rank \(=\) row rank \(=p\), so \(A\) has a right inverse \(Z\in\Delta^{q\times p}\), \(A(Zt)=t\) for all \(t\in\Delta_{\mathbb R}^p\). A continuous surjective homomorphism of compact groups maps normalised Haar measure to normalised Haar measure (the image measure is a translation invariant probability measure). Hence \(\int_{\mathcal P}\tilde\psi(An)dn=\int_{\mathbb T'}\tilde\psi\). \(\square\) (No primitivity of \(A\) is needed; full rank is.)

\textbf{Lemma 2.4 (fibre mean value at \(e_1\)) {[}P from (R\(_{j-1,k-1}\)){]}.} Let \(2\le j<k\), assume (R\(_{j-1,k-1}\)) and that \(\Gamma_{k-1}\) is a lattice in \(\mathsf G_{k-1}\). For every Borel \(\Psi:\Delta_{\mathbb R}^{(j-1)\times k}\to[0,\infty]\):
\[\int_{\Gamma_{\mathsf H}\backslash\mathsf H}\ \sum_{Y'\in\mathcal O^{(j-1)\times k},\ \mathrm{rk}\binom{e_1}{Y'}=j}\Psi(Y'h)\ d\eta(h)=\int_{\Delta_{\mathbb R}^{(j-1)\times k}}\Psi(Z)\,dZ .\]
\emph{Proof.} Write \(Y'=(b\,|\,Y'')\), \(b\in\mathcal O^{(j-1)\times1}\), \(Y''\in\mathcal O^{(j-1)\times(k-1)}\). Row operations give \(\mathrm{rk}\binom{e_1}{Y'}=1+\mathrm{rk}\,Y''\), and \(Y'u(n,h)=(b+Y''n\,|\,Y''h)\). The integrand is left \(\Gamma_{\mathsf H}\)-invariant (\(Y'\mapsto Y'\epsilon\) permutes the index set for \(\epsilon\in\Gamma_{\mathsf H}\)). By Lemma 2.2(iv), Tonelli, and Lemma 2.3 with \(A=Y''\), \(p=j-1\), \(q=k-1\):
\[\text{LHS}=\int_{\mathcal F'}\sum_{\mathrm{rk}Y''=j-1}\Big[\int_{\mathcal P}\sum_b\Psi(b+Y''n\,|\,Y''h)\,dn\Big]dh=\int_{\Gamma_{k-1}\backslash\mathsf G_{k-1}}\sum_{\mathrm{rk}Y''=j-1}\Psi_1(Y''h)\,dh,\]
\(\Psi_1(Z''):=\int\Psi(t\,|\,Z'')dt\) (Borel, \([0,\infty]\)-valued). By (R\(_{j-1,k-1}\)) this is \(\int\Psi_1(Z'')dZ''=\int\Psi\). \(\square\)

\textbf{Lemma 2.5 (fibre mean value at any \(y\)) {[}P{]}.} Same hypotheses and (h1\(^\ell\)). Let \(0\ne y\in\mathcal O^{1\times k}\), \(\mathsf H_y:=\{g\in\mathsf G_k:yg=y\}\), \(\Gamma_y:=\mathsf H_y\cap\Gamma_k\). Then \(\mathsf H_y\) is unimodular, \(\Gamma_y\) is a lattice in it, and with \(\eta_y\) the Haar measure of \(\mathsf H_y\) normalised by \(\eta_y(\Gamma_y\backslash\mathsf H_y)=1\):
\[\int_{\Gamma_y\backslash\mathsf H_y}\sum_{Y_1\in\mathcal O^{(j-1)\times k},\ \mathrm{rk}\binom{y}{Y_1}=j}\Psi(Y_1h)\,d\eta_y(h)=\int\Psi(Z)\,dZ\qquad(\Psi\ge0\text{ Borel}).\]
\emph{Proof.} \(y=ae_1\delta\) (Lemma 2.1). \(a\) is invertible in \(\Delta_{\mathbb R}\) and \(\delta\) normalises \(\mathsf G_k\) (G1) and \(GL_k(\mathcal O)\), hence \(\Gamma_k\). So \(\mathsf H_y=\delta^{-1}\mathsf H\delta\), \(\Gamma_y=\delta^{-1}\Gamma_{\mathsf H}\delta\), and \(h_1\mapsto\delta^{-1}h_1\delta\) carries \(\eta\) to \(\eta_y\). Substituting \(h=\delta^{-1}h_1\delta\), \(Y_1=Y'\delta\) (a bijection of \(\mathcal O^{(j-1)\times k}\); \(\mathrm{rk}\binom{y}{Y_1}=\mathrm{rk}\binom{ae_1}{Y'}=\mathrm{rk}\binom{e_1}{Y'}\)), \(Y_1h=Y'h_1\delta\), the left side becomes the left side of Lemma 2.4 for \(\Psi^\delta(Z):=\Psi(Z\delta)\), and \(\int\Psi^\delta=\int\Psi\) by Lemma 0.2. \(\square\)
(The non-primitive part \(a\) does not enter the fibre integral at all.)

\textbf{Lemma 2.6 (unfolding and disintegration) {[}L, classical (Weil\textquotesingle s quotient integral formula; e.g. Raghunathan, \emph{Discrete subgroups of Lie groups}, Ch. I), from memory + P{]}.} Let \(\mathsf G\) be a unimodular lcsc group, \(\Gamma\) a lattice, \(m\) the Haar measure with \(m(\Gamma\backslash\mathsf G)=1\), \(\mathsf H'\subseteq\mathsf G\) a closed unimodular subgroup such that \(\Gamma':=\Gamma\cap\mathsf H'\) is a lattice in \(\mathsf H'\), \(\eta'\) the Haar measure of \(\mathsf H'\) with \(\eta'(\Gamma'\backslash\mathsf H')=1\). Then there is a \(\mathsf G\)-invariant measure \(\nu'\) on \(\mathsf H'\backslash\mathsf G\) such that for every left-\(\Gamma'\)-invariant Borel \(\Phi:\mathsf G\to[0,\infty]\)
\[\int_{\Gamma\backslash\mathsf G}\sum_{\delta\in\Gamma'\backslash\Gamma}\Phi(\delta g)\,dm=\int_{\Gamma'\backslash\mathsf G}\Phi\,dm=\int_{\mathsf H'\backslash\mathsf G}\Big[\int_{\Gamma'\backslash\mathsf H'}\Phi(hg)\,d\eta'(h)\Big]d\nu'(\mathsf H'g).\]
\emph{Proof.} First equality: if \(\mathcal F\) is a Borel fundamental domain for \(\Gamma\), \(\bigsqcup_{\delta\in\Gamma'\backslash\Gamma}\delta\mathcal F\) is one for \(\Gamma'\) (Tonelli). Second: Weil\textquotesingle s formula \(\int_{\mathsf G}\phi\,dm=\int_{\mathsf H'\backslash\mathsf G}\int_{\mathsf H'}\phi(hg)d\eta'(h)d\nu'\) (valid for Borel \(\phi\ge0\); existence of \(\nu'\) because both groups are unimodular) applied to \(\phi=\Phi1_{\mathcal E}\), \(\mathcal E\) a fundamental domain for \(\Gamma'\) in \(\mathsf G\): with \(\mathcal D\) a fundamental domain of \(\Gamma'\) in \(\mathsf H'\), \(\int_{\mathsf H'}\Phi(hg)1_{\mathcal E}(hg)d\eta'=\int_{\mathcal D}\Phi(hg)\sum_{\epsilon\in\Gamma'}1_{\mathcal E}(\epsilon hg)d\eta'=\int_{\mathcal D}\Phi(hg)d\eta'\). \(\square\)

\hypertarget{proof-of-theorem-r-induction-step-jge2-kge-j1ge3}{%
\subsubsection*{\texorpdfstring{3. Proof of Theorem R (induction step \(j\ge2\), \(k\ge j+1\ge3\))}{3. Proof of Theorem R (induction step j\textbackslash ge2, k\textbackslash ge j+1\textbackslash ge3)}}\label{proof-of-theorem-r-induction-step-jge2-kge-j1ge3}}

Let \(m_k\) be the Haar measure of \(\mathsf G_k\) inducing \(\mu_k\). \(\mathcal O^{1\times k}\setminus\{0\}\) is a countable disjoint union of \(\Gamma_k\)-orbits \(y_i\Gamma_k\), \(i\in I\) (NO transitivity is claimed or needed: \(I\) is infinite, it contains all non-primitive classes and possibly several primitive classes). Put \(\mathsf H_i=\mathsf H_{y_i}\), \(\Gamma_i=\Gamma_{y_i}\), \(\eta_i=\eta_{y_i}\) (Lemma 2.5), \(\nu_i\) the measure of Lemma 2.6 for \((\mathsf G_k,\Gamma_k,\mathsf H_i)\).

\emph{Step 1 (orbit decomposition).} \(\delta\mapsto y_i\delta\) is a bijection \(\Gamma_i\backslash\Gamma_k\to y_i\Gamma_k\). Writing \(Y=\binom{y}{Y'}\), \(y=y_i\delta\), \(Y'=Y_1\delta\) (\(\mathrm{rk}\binom{y_i\delta}{Y_1\delta}=\mathrm{rk}\binom{y_i}{Y_1}\)):
\[S_jF(\gamma):=\sum_{\mathrm{rk}Y=j}F(Y\gamma)=\sum_{i\in I}\sum_{\delta\in\Gamma_i\backslash\Gamma_k}\Phi_i(\delta\gamma),\qquad\Phi_i(g):=\sum_{Y_1\in\mathcal O^{(j-1)\times k},\ \mathrm{rk}\binom{y_i}{Y_1}=j}F\binom{y_ig}{Y_1g}.\]
\(\Phi_i\) is Borel, \([0,\infty]\)-valued and left \(\Gamma_i\)-invariant (\(Y_1\mapsto Y_1\epsilon\) permutes the index set for \(\epsilon\in\Gamma_i\)).

\emph{Step 2 (unfold, disintegrate).} By Tonelli and Lemma 2.6: \(\int_{X_k}S_jF\,d\mu_k=\sum_i\int_{\mathsf H_i\backslash\mathsf G_k}\big[\int_{\Gamma_i\backslash\mathsf H_i}\Phi_i(hg)d\eta_i(h)\big]d\nu_i(g)\).

\emph{Step 3 (fibre).} For fixed \(g\), since \(y_ih=y_i\): \(\Phi_i(hg)=\sum_{Y_1}\Psi_g(Y_1h)\) with \(\Psi_g(Z):=F\binom{y_ig}{Zg}\). By Lemma 2.5 and \(\|g\|=1\): \(\int_{\Gamma_i\backslash\mathsf H_i}\Phi_i(hg)d\eta_i(h)=\int F\binom{y_ig}{Zg}dZ=\int F\binom{y_ig}{Z}dZ=f(y_ig)\), where \(f(x):=\int_{\Delta_{\mathbb R}^{(j-1)\times k}}F\binom{x}{Z}dZ\) is Borel \(\Delta_{\mathbb R}^{1\times k}\to[0,\infty]\). Hence
\[\int_{X_k}S_jF\,d\mu_k=\sum_{i\in I}\int_{\mathsf H_i\backslash\mathsf G_k}f(y_ig)\,d\nu_i(g).\tag{3.1}\]

\emph{Step 4 (the constant, via (L1)\(_k\); point (d)).} Run Steps 1-2 for \(j=1\) with \(f\) in place of \(F\): \(\Phi_i(g)=f(y_ig)\) is left \(\mathsf H_i\)-invariant, the fibre integral is \(\eta_i(\Gamma_i\backslash\mathsf H_i)f(y_ig)=f(y_ig)\), so \(\int_{X_k}\sum_{y\ne0}f(y\gamma)d\mu_k=\sum_i\int_{\mathsf H_i\backslash\mathsf G_k}f(y_ig)d\nu_i(g)\): the SAME right side as (3.1), orbit by orbit, with the same measures \(\nu_i\). By (L1)\(_k\) and Lemma 0.5 the left side is \(\int f(x)dx=\int F\) (Tonelli). So \(\int_{X_k}S_jF\,d\mu_k=\int F\). \(\blacksquare\)

\textbf{Remarks.} (1) \emph{Points (b),(c) of the task.} The measures \(\nu_i\) on \(\mathsf H_i\backslash\mathsf G_k\cong y_i\mathsf G_k\) (for families containing \(\mathsf G^1_k\), \(k\ge2\): the open dense orbit \(\{x:x_v\ne0\ \forall v\}\)) are never identified with multiples of Lebesgue measure, and the number and the volumes of the orbits are never computed: (L1)\(_k\) says exactly that \(\sum_i\nu_i\)-integrals reproduce \(\int f\), and the fibre mean value (Lemma 2.5) is the same number 1 for every orbit, primitive or not. This is why no zeta value appears. Classical check (\(\Delta=\mathbb Q\)): \(\nu\)-contributions are \(\zeta(k)^{-1}a^{-k}\int f\) for \(y=ae_1\), summing to \(\int f\). (2) \emph{Sanity check \(\Delta=\mathbb Q\), \(\mathcal O=\mathbb Z\):} Theorem R is the principal term of Rogers\textquotesingle{} formula: for \(j\le k-1\), \(\mathbb E\sum_{\text{lin. indep. }(y_1..y_j)\in(L)^j}F=\int F\) {[}L, Rogers 1955, from memory{]}; \(j=k\) is false classically, and here the induction would need (L1) in rank 1, which fails. (3) \(\Delta_{\mathbb R}\cong\mathbb K^s\) is used only in Lemma 0.2/0.3; for general \(\Delta_{\mathbb R}\) the same proof works once (G1), (G2) are checked {[}not done here{]}. (4) Integrability is never needed: all identities are in \([0,\infty]\) by Tonelli; for signed \(F\) apply to \(F^\pm\) when \(\int|F|<\infty\). (5) \emph{Transitivity (point (b)), for information only.} By Lemma 2.1, \(GL_k(\mathcal O)\) is transitive on primitive vectors {[}P{]}. For the family \(\mathsf G^\pm\), \(\Gamma_k=GL_k(\mathcal O)\) (Lemma 0.2), so there is one primitive orbit {[}P{]}. For \(\mathsf G^+\) or family (A), \(\Gamma_k=\ker(\chi_k|GL_k(\mathcal O))\) may be smaller and the primitive vectors split into \([\chi_k(GL_k(\mathcal O)):\chi_{k-1}(GL_{k-1}(\mathcal O))]\) orbits (the \(GL_k(\mathcal O)\)-stabiliser of \(e_1\) is \(\{u(n,h):h\in GL_{k-1}(\mathcal O)\}\)); whether this index is 1 is a \(K_1\)-stability question that I did NOT settle and that the proof does not need: Step 4 matches every orbit separately. (6) \emph{Without (h1\(^\ell\))} the saturation \(\Delta y\cap\mathcal O^{1\times k}\) need not be a free direct summand with free complement; the fibre over \(y\) is then a mean value problem for a non-free \(\mathcal O\)-lattice of rank \(k-1\) and its stabiliser, which (L1) as cited (free lattices only) does not cover \textbf{{[}GAP if (h1\(^\ell\)) is dropped{]}}. (7) \emph{Where \(\mathrm{rk}\,Y=j\) and \(j<k\) are used:} full rank of \(Y''\) makes the torus map of Lemma 2.3 onto (for lower rank the \(b\)-sum does not become a full Lebesgue integral: these are the lower-rank terms of Rogers\textquotesingle{} formula, which {[}R14, Lemma 1.2{]} treats separately); \(j<k\) keeps every rank \(\ge2\).

\hypertarget{status-gaps-what-a-referee-should-check}{%
\subsubsection*{4. Status, gaps, what a referee should check}\label{status-gaps-what-a-referee-should-check}}

\textbf{4.1 Left/right form of (h1) {[}flag{]}.} The proof needs \textbf{(h1\(^\ell\))} (left ideals), because rows form a LEFT \(\mathcal O\)-module; {[}R14{]} states (h1) for RIGHT ideals (needed there for \(\Lambda_W\)). If \(\mathcal O\) is stable under an anti-automorphism of \(\Delta\) (commutative \(\Delta\); a quaternion algebra with \(\mathcal O\) stable under the canonical involution, e.g. \(\mathcal O\ni\mathrm{trd}(x)\) for \(x\in\mathcal O\): Hurwitz, icosians, and the Leech order if it is an \(\mathcal O_F\)-order) the two forms are equivalent {[}P, trivial{]}. For maximal orders they are equivalent via \(I\mapsto I^{-1}\) {[}L, from memory{]}; in general via duality for hereditary noetherian rings {[}L, from memory, not checked{]}. The task statement phrases (h1) as "every one-sided ideal is principal", which covers both.

\textbf{Lemma 4.1a (left versus right in the setting (T) of {[}R14{]}) {[}P{]}.} Let \(B\) be a rational \(G\)-invariant positive definite symmetric form on \(W\), \(\Lambda_W^\#\) the \(B\)-dual lattice. Then there is an anti-automorphism \(c\mapsto c^*\) of \(\Delta\) with \(B(xc,y)=B(x,yc^*)\), and \(\mathcal O(\Lambda_W^\#)=\mathcal O(\Lambda_W)^*\). Hence (h1\(^\ell\)) for \(\mathcal O=\mathcal O(\Lambda_W)\) \(\iff\) every nonzero right ideal of \(\mathcal O(\Lambda_W^\#)\) is principal; in particular \textbf{if \(\Lambda_W\) is \(B\)-unimodular (Leech lattice with its standard form, which is \(G\)-invariant) then (h1) \(\iff\) (h1\(^\ell\)).}
\emph{Proof.} The \(B\)-adjoint of \(R_c:x\mapsto xc\) is \(\mathbb Q\)-linear and commutes with \(G\) (\(G\) acts by \(B\)-isometries commuting with \(R_c\)), so it is \(R_{c^*}\) for a unique \(c^*\in\Delta=\mathrm{End}_G(W)^{op}\); \(R_{cc'}=R_{c'}R_c\) gives \((cc')^*=c'^*c^*\), and \(c^{**}=c\) since \(B\) is symmetric. If \(\Lambda_Wc\subseteq\Lambda_W\) and \(y\in\Lambda_W^\#\) then \(B(x,yc^*)=B(xc,y)\in\mathbb Z\) for \(x\in\Lambda_W\), so \(\Lambda_W^\#c^*\subseteq\Lambda_W^\#\); the converse by \(\Lambda_W^{\#\#}=\Lambda_W\). \(J\mapsto J^*\) exchanges left ideals of \(\mathcal O\) and right ideals of \(\mathcal O^*\), principal ones with principal ones. \(\square\)
So for Theorem D (Leech) of {[}R14{]} no new hypothesis is needed; for Theorem E in general, (h1\(^\ell\)) is an additional hypothesis unless \(\mathcal O(\Lambda_W^\#)=\mathcal O(\Lambda_W)\) (or one of the equivalences above applies). \textbf{A referee should check Lemma 4.1a and that the Leech form used in {[}R14{]} is the standard unimodular one.}

\textbf{4.2 Which group is in {[}G{]}? {[}flag, not a gap of the proof{]}.} Theorem R is proved for every coherent family satisfying (L1) in the ranks used. {[}R6, Rem. 4.1a(i){]} records \(\mathsf G=\{N=1\}\) for {[}G{]}; the task text says "reduced norm one at every place". If no \(\mathbb K_v=\mathbb R\), \(\{N=1\}=\mathsf G^+_m=\mathsf G^\pm_m\) is coherent and Theorem R applies verbatim to {[}R14{]} (Leech, Hurwitz, icosians). If some \(\mathbb K_v=\mathbb R\) and \(N\) is the determinant on \(M_m(\Delta_{\mathbb R})\) (an \(m\)-th power), \(\{N=1\}\) in rank \(m\) and the Levi factor of the stabiliser in rank \(m+1\) may differ by an index-2 parity; then one passes to the coherent family \(\mathsf G^\pm\) or \(\mathsf G^+\) by the following lemma (written last, concise, \textbf{un-refereed; only needed when a real place exists}). (L1) for family (A) implies (L1) for family (B) by averaging over the compact quotient, not conversely {[}H, not proved here{]}.

\textbf{Lemma 4.2a (finite-index change of group) {[}P, concise{]}.} Let \(m\ge2\) and \(\mathsf G'\subseteq\mathsf G''\) be closed subgroups of \(\{\|\cdot\|=1\}\subseteq GL_m(\Delta_{\mathbb R})\), \(\mathsf G'\) normal of finite index in \(\mathsf G''\), \(\mathsf G'\supseteq\mathsf G^1_m\); \(\Gamma'=\mathsf G'\cap GL_m(\mathcal O)\), \(\Gamma''=\mathsf G''\cap GL_m(\mathcal O)\). Then (L1)\(_m\) for \((\mathsf G',\Gamma')\) \(\iff\) (L1)\(_m\) for \((\mathsf G'',\Gamma'')\).
\emph{Proof.} \(\Gamma'\) has finite index in \(\Gamma''\) and \(\mathsf G'\) is open of finite index in \(\mathsf G''\), so the lattice properties are equivalent. \(\Gamma''\mathsf G'\) is a subgroup of finite index \(L\); \(\mathsf G''=\bigsqcup_{l\le L}\Gamma''\mathsf G'\kappa_l\), \(\kappa_1=1\), and \(\Gamma'g\mapsto\Gamma''g\kappa_l\) is a bijection of \(\Gamma'\backslash\mathsf G'\) onto the \(l\)-th piece of \(\Gamma''\backslash\mathsf G''\) carrying \(\mu'\) to \(L\mu''|_{\text{piece}}\) (invariance under \(\mathsf G'=\kappa_l^{-1}\mathsf G'\kappa_l\), uniqueness). Hence with \(\lambda'(A)=\mathbb E'\#\{y\ne0:y\gamma\in A\}\), \(\lambda''\) likewise: \(\lambda''=\frac1L\sum_l(R_{\kappa_l})_*\lambda'\). "\(\Rightarrow\)": \(\lambda'=\) Leb and \(\|\kappa_l\|=1\). "\(\Leftarrow\)": \(\lambda'\le L\cdot\)Leb, so \(\lambda'\) is a \(\mathsf G'\)-invariant Radon measure vanishing on the null set \(\{x:\exists v,\ x_v=0\}\); on the complement \(\Omega\), \(\mathsf G'\supseteq\prod_vSL_m(\mathbb K_v)\) is transitive (\(m\ge2\)) with unimodular stabiliser, and Leb\(|_\Omega\) is invariant, so \(\lambda'=c\,\)Leb (uniqueness of the invariant measure on a homogeneous space {[}L, classical{]}); then \(c=1\) from the average. \(\square\)

\textbf{4.3 Cited, not proved.} (L1)\(_m\) including "\(\Gamma_m\) is a lattice" ({[}G{]}, not opened in this session; caveat of {[}R6, Rem. 4.1a(iv){]} on the proof of {[}G, Thm 16{]} for \(s\ge2\) stands). Weil\textquotesingle s quotient integral formula, existence of Borel fundamental domains, "a group with a lattice is unimodular", Haar measure under surjective homomorphisms of compact groups, row rank = column rank over a division ring, generation of \(GL_m\) by elementary and diagonal matrices: classical, from memory.

\textbf{4.4 Things to check.} (i) Lemma 2.1: the left-handed version of {[}R14, Lemma 1.1{]} and that \(\delta^{-1}\) is integral. (ii) Lemma 2.2(ii),(iv): Haar measure and fundamental domain of \(\mathsf H\); uses \(\|h\|=1\) on COLUMNS (Lemma 0.2). (iii) Lemma 2.3: surjectivity needs full rank of \(Y''\), which is exactly the condition \(\mathrm{rk}\,Y=j\). (iv) Step 4: the same \(\nu_i\) (normalised through \(\eta_i(\Gamma_i\backslash\mathsf H_i)=1\)) appear for \(j\) and for \(1\). (v) Measurability of \(f\) and \(\Psi_1\) with values in \([0,\infty]\) (Tonelli). (vi) The set of ranks: (R\(_{j,k}\)) uses (L1) in ranks \(k,\dots,k-j+1\); for {[}R14{]} (\(j\le r<k\)) all are \(\ge2\).

\textbf{4.5 Outcome for R14.} (H-Rogers)\(_j\), \(2\le j\le r<k\), is no longer a hypothesis: it follows from (L1) (in ranks \(k-r+1,\dots,k\)) and (h1\(^\ell\)), with constant exactly 1; Theorems D/E of {[}R14{]} then depend only on (L1)-(L3), (BS), the R5 facts and (h1), (h1\(^\ell\)).
